% BEGIN R28_BOUNDARY_OCCUPANCY
% Insert in MG01 after the definition of c_j,o_j, before their values.
The law for the boundary at depth six fixes \(c_j=1\). With the axis \(z\) and the
aggregate \(q\) already determined, the integer lateral sums are
\[
 (x_j+y_j)_{j=1}^6=(2,4,3,4,4,2).
\]
The minimum-occupancy rule selects, among decompositions with
\(x_j,y_j\in\{0,1,2\}\), the columns with the fewest nonzero entries,
including \(z_j\). For sum two, the extreme pairs \((0,2),(2,0)\)
occupy one lateral position; for sum three, the pairs \((1,2),(2,1)\)
occupy two; for sum four, only \((2,2)\) is possible. Including the
axis gives \(o=(2,2,3,2,3,2)\). At this boundary, with
\(s=[x+y]_3\), the same rule is expressed by
\[
 o_j=2+\mathbf1_{\{z_j\ne0,\ s_j\ne2\}}.
\]
% END R28_BOUNDARY_OCCUPANCY
